\documentclass[10pt,fleqn,a4paper]{jsarticle}
\usepackage[dvipdfmx]{graphicx,color}
\usepackage{ascmac,amsmath,amssymb,amstext}
\usepackage{tikz,multicol,float,tikz-3dplot}
\usetikzlibrary{positioning,intersections,calc,arrows.meta,math,angles}
\usepackage{zogeny,ceo,comment}
\newcommand{\ans}[1]{\mbox{\boldmath{$#1$}}}
\renewcommand{\baselinestretch}{1.2} 
\tdplotsetmaincoords{60}{110}%{xy平面からどれだけ上にいるか(90度から-90度)}{z軸中心にどれだけ左右に回転するか}
\title{二項定理}
\author{大島　遙斗}
\begin{document}
\maketitle
\newpage
割と皆さんが苦手な二項定理です.ですが,二項定理が苦手な理由は,教科書や参考書に騙されているからである可能性が高いです.ですので,
全ての二項定理の公式は\ans{あたりまえ}にしてしまいましょう.\\
\begin{itembox}
[l]{\enpitu 定理}
次のような展開を二項展開とよび,数学的帰納法で証明することができる.
\begin{center}
    $\ans{(x+y)^n=\wa{k=0}{n}\comb{n}{k}x^ky^{n-k}=\wa{k=0}{n}\comb{n}{k}x^{n-k}y^k}$
\end{center}
\end{itembox}\\
\begin{itembox}
    [l]{\enpitu\hspace{2mm}$\comb{n}{k}の性質$}
    \doichi. $\comb{n}{k}=\comb{n}{n-k}　(0\leq k\leq n)\\
    \doni.\comb{n}{k}=\comb{n-1}{k-1}+\comb{n-1}{k}　(1\leq k\leq n-1)\\
    \dosan. k\cdot\comb{n}{k}=n\cdot\comb{n-1}{k-1}　(1\leq k\leq n)$
\end{itembox}
\newpage
\noindent\ba{1}.次のそれぞれの計算を実行せよ.\\
\kakkoichi$(x+y)^7を展開したときのx^3y^4の係数を求めよ.\\
\kakkoni (x+y+z)^8を展開したときのx^2yz^5の係数を求めよ.$\\\\
\ba{2}.次の和を求めよ.ただし,$nは自然数である.$\\
\kakkoichi $A_n=\wa{k=0}{n}3^k\cdot\comb{n}{k}\\
\kakkoni B_n=\wa{k=0}{n}\comb{n}{k}\\
\kakkosan C_n=\wa{k=1}{n}\comb{2n}{2k-1}$\\\\

\noindent\fbox{\kurosankakub 日々の演習\kurosankakua}\\
\shikakuichib.次の多項式を展開したときの指定された文字の係数を求めよ.\\
\kakkoichi $(2x-1)^6のx^4の係数.\\
\kakkoni (x+2y+3z)^7のx^2y^3z^2の係数.$\\
\shikakunib.次の等式を証明せよ.
\begin{align*}
    \wa{k=0}{n}\p{\comb{n}{k}}^2=\comb{2n}{n}
\end{align*}
\shikakusanb.次の和を求めよ.\\
$\kakkoichi \wa{k=0}{n}\dfrac{\comb{n}{k}}{k+1}\\
\kakkoni \wa{k=0}{n}\dfrac{\comb{2n}{2k}}{2k+1}$\\\\
\shikakushib.$\hspace{1mm}3^{100}の下二桁を求めよ.\hspace{1cm}(お茶の水女子大)$\\\\
\shikakugob.次の各問いに求めよ.\\
$\kakkoichi nを正の整数とする.(1+a)^nを二項定理を用いて展開せよ.\\
\kakkoni 21^{10}を400で割った余りを求めよ.\\
\kakkosan 19^n+21^nが400で割り切れるような正の整数nは存在するか.存在するならばその例を示し,存在しない場合はそれを証明せよ.\hspace{1cm}(中央大)$
\newpage
{\large \ans{解答}}\\
\ba{1}.\\
\kai\\
\kakkoichib $xを3個,yを4個選べばよいので,求める係数は,\comb{7}{3}=\dfrac{7\cdot6\cdot5}{3\cdot2\cdot1}=\ans{35}\\
\kakkonib xを2個,yを1個,zを5個選べば良いので,\comb{8}{2}x^2\cdot\comb{6}{1}y\cdot\comb{5}{5}z^5=\dfrac{8\cdot7}{2\cdot1}\times\dfrac{6}{1}\times1x^2yz^5=\ans{168}$\\
\bekkai\kakkoni\\
$xを2個,yを1個,zを5個並べるときの順列の総数に等しいので,\dfrac{8!}{2!5!}=\dfrac{8\cdot7\cdot6}{2\cdot1}=\ans{168}$\\\\
\ba{2}.\\
\kangaekata\\
\kakkoichi \kakkoni どうにかして,二項定理の定義の形に帰着させることを考えましょう.\\
\kakkosan よくわからない数列は,\ans{書き下す}が基本です.\\
\kai\\
\kakkoichib $A_n=\wa{k=0}{n}3^k\cdot\comb{n}{k}=\wa{k=0}{n}\comb{n}{k}\cdot3^k\cdot1^{n-k}=(3+1)^n=\ans{4^n}\\
\kakkonib B_n=\wa{k=0}{n}\comb{n}{k}=\wa{k=1}{n}\comb{n}{k}1^{k}\cdot1^{n-k}=(1+1)^n=\ans{2^n}$\\
\kakkosanb $C_n=\wa{k=1}{n}\comb{2n}{2k-1}=\comb{2n}{1}+\comb{2n}{3}+\comb{2n}{5}+\cdots+\comb{2n}{2n-1}\\
これは,(1+1)^{2n}を展開したときの,奇数番目のみ抜き出したものである.\\
このことより,$
\begin{align*}
    (1+1)^{2n}&=\comb{2n}{0}+\comb{2n}{1}+\comb{2n}{2}+\comb{2n}{3}+\cdots+\comb{2n}{2n}\cdots\cdots\cdots\maruichi\\
    (-1+1)^{2n}&=\comb{2n}{0}-\comb{2n}{1}+\comb{2n}{2}-\comb{2n}{3}+\cdots+\comb{2n}{2n}\cdots\cdots\cdots\maruni
\end{align*}
$よって,\dfrac{\maruichi-\maruni}{2}より,$
\begin{align*}
    \dfrac{\maruichi-\maruni}{2}=\dfrac{2^{2n}-0}{2}=\comb{2n}{1}+\comb{2n}{3}+\comb{2n}{5}+\cdots+\comb{2n}{2n-1}
\end{align*}
なので,求めたい和と一致し,$C_n=\ans{\dfrac{2^{2n}}{2}}=\ans{2^{2n-1}}$\\\\
\newpage
\noindent{\large\ans{\fbox{\kurosankakub 日々の演習\kurosankakua}の解答}}\\
\noindent\shikakuichib\\
\noindent\kai\\
\kakkoichib $2xを4個選べば良いので,求める係数は,\comb{6}{4}(2x)^4(-1)^2=15\cdot16=\ans{240}$\\
\kakkonib $xを2個,2yを3個,3zを2個選べば良いので,\comb{7}{2}\times\comb{5}{3}2^3\times3^2=\ans{15120}$\\\\
\shikakunib\\
\kangaekata\\
帰納法は明らかにめんどくさそうです.なので,一回直接証明することができないかやってみましょう.\\
\kai\\
$\wa{k=0}{n}\p{\comb{n}{k}}^2=\comb{2n}{n}\cdots\cdots\cdots\asta　を示す.$
\begin{align*}
    (1+x)^{2n}=(1+x)^n\times(1+x)^n
\end{align*}
より,二項定理から
\begin{align*}
    \comb{2n}{0}+\comb{2n}{1}x+\comb{2n}{2}x^2+\cdots+\comb{2n}{n}x^n+\cdots+\comb{2n}{2n}x^{2n}&=\p{\comb{n}{0}+\comb{n}{1}x+\comb{n}{2}x^2+\cdots+\comb{n}{n}x^n}\\
    &\times\p{\comb{n}{0}+\comb{n}{1}x+\comb{n}{2}x^2+\cdots+\comb{n}{n}x^n}\cdots\cdots\cdots\maruichi
\end{align*}
$\maruichi の両辺のx^nの係数を比較すると$
\begin{align*}
    \comb{2n}{n}&=\comb{n}{0}\times\comb{n}{n}+\comb{n}{1}\times\comb{n}{n-1}+\comb{n}{2}\times\comb{n}{n-2}+\cdots+\comb{n}{n}\times\comb{n}{0}\\
    &=\wa{k=0}{n}=\comb{n}{k}\times\comb{n}{n-k}\\
    &=\wa{k=0}{n}\comb{n}{k}\times\comb{n}{k}　\p{\because \comb{n}{n-k}=\comb{n}{k}}\\
    &=\wa{k=0}{n}\p{\comb{n}{k}}^2
\end{align*}
よって,$ \asta が示された.\ans{(証明終わり)}$\\\\
\newpage
\noindent\shikakusanb\\
\kangaekata\\
\kakkoichi\kakkoni 分母に式が来るような計算は...と考えると\ans{積分すれば良い}と思い付きます.\\
\kai\\
\kakkoichib 二項定理より,$(1+x)^n=\wa{k=0}{n}\comb{n}{k}x^kなので,これの両辺を0から1で積分すると,$
\begin{align*}
    &\dint{0}{1}(1+x)^ndx=\dint{0}{1}\wa{k=0}{n}\comb{n}{k}x^kdx\\
    &\doti \tint{\dfrac{(1+x)^{n+1}}{n+1}}{0}{1}=\tint{\wa{k=0}{n}\dfrac{\comb{n}{k}}{k+1}x^{k+1}}{0}{1}\\
    &\doti \dfrac{2^{n+1}-1}{n+1}=\wa{k=0}{n}\dfrac{\comb{n}{k}}{k+1}
\end{align*}
よって,これは求める和と一致し,求める和は$\ans{\dfrac{2^{n+1}-1}{n+1}}$\\
\kakkonib 二項定理より,$(1+x)^{2n}=\wa{k=0}{2n}\comb{2n}{k}x^k=\underbrace{\wa{k=0}{n}\comb{2n}{2k}x^{2k}}_{偶数番目}+\underbrace{\wa{k=1}{n}\comb{2n}{2k-1}x^{2k-1}}_{奇数番目}
(\temarkua 数\suIII の無限級数でもよくやるテクニックです.)\\$
これの両辺を$-1から1まで積分すると,$
\begin{align*}
    &\dint{-1}{1}(1+x)^{2n}dx=2\dint{0}{1}\wa{k=0}{n}\comb{2n}{2k}x^{2k}dx\\
    &\doti \tint{\dfrac{(1+x)^{2n+1}}{2n+1}}{-1}{1}=2\tint{\wa{k=0}{n}\dfrac{\comb{2n}{2k}}{2k+1}}{0}{1}\\
    &\doti \dfrac{2^{2n+1}}{2n+1}=2\wa{k=0}{n}\dfrac{\comb{2n}{2k}}{2k+1}\\
    &\doti \wa{k=0}{n}\dfrac{\comb{2n}{2k}}{2k+1}=\dfrac{4^n}{2n+1}
\end{align*}
よって,求める和は,$\ans{\dfrac{4^n}{2n+1}}$\\\\
\newpage
\noindent\shikakugob\\
\kangaekata\\
$下二桁なので\ans{100で割った余り}ですね!二項定理で展開して100で割り切れないところだけ残しましょう.$\\
\kai\\
二項定理より,
\begin{align*}
    3^{100}&=(1+2)^{100}\\
    &=\wa{k=0}{100}\comb{100}{k}2^k\cdots\cdots\cdots\maruichi
\end{align*}
$\maruichi の右辺を100で割ると,\comb{100}{0}のみ残るので,求める余りは,\ans{1}$\\\\
\shikakugob\\
\kangaekata\\
\kakkosan 二項定理で数を小さくしましょう.\\
\kai\\
\kakkoichib $(1+a)^n=\ans{\wa{k=0}{n}\comb{n}{k}a^k}$ \\
\kakkonib 
\begin{align*}
    21^{10}&=(20+1)^{10}=\wa{k=0}{10}\comb{10}{k}20^k\\
    &=\comb{10}{0}+\comb{10}{1}20+\wa{k=2}{10}\comb{10}{k}20^k\cdots\cdots\cdots\maruichi
\end{align*}
ここで,$\maruichi の左辺において,\wa{k=2}{10}\comb{10}{k}20^k は,必ず400を因数に持つので400で割り切れる.$\\
$\Y 21^{10}を400で割った余りは,\comb{10}{0}+\comb{10}{1}20=\ans{201}$\\\\
\kakkosanb $f(n)=19^n+21^nとおく.$
\begin{align*}
    f(n)が400で割り切れるような正の整数nが存在する.\doti f(n)\godo0\kmod{400}\cdots\cdots\cdots\asta
\end{align*}
これより,$19^nと21^nをそれぞれ二項定理で展開すると$
\begin{align*}
    19^n&=(20-1)^n\\
    &=\wa{k=0}{n}\comb{n}{k}\p{20}^k\p{-1}^{n-k}\\
    &=\comb{n}{0}(-1)^n+\comb{n}{1}20\p{-1}^{n-1}+\wa{k=2}{n}\comb{n}{k}\p{20}^k\p{-1}^{n-k}\\
    &\godo \comb{n}{0}(-1)^n+\comb{n}{1}20\p{-1}^{n-1}\kmod{400}\\
    &=\p{-1}^n\p{1-20n}\kmod{400}\cdots\cdots\cdots\maruni
\end{align*}
\begin{align*}
    21^n&=\comb{n}{0}+\comb{n}{1}20+\wa{k=2}{n}\comb{n}{k}20^k\\
    &\godo1+20n\kmod{400}\cdots\cdots\cdots\marusan
\end{align*}
よって,$\maruni,\marusan から,f(n)\godo \p{-1}^n\p{1-20n}+1+20n$\\
\tokeiichi\underline{$nが偶数のとき,$}\\
　$(-1)^n=1なので,f(n)=\p{1-20n}+\p{1+20n}=2\\
これより,f(n)\negodo 0\kmod{400}.$\\
\tokeini\underline{$nが奇数のとき,$}\\
　$(-1)^n=-1なので,f(n)=(-1)(1-20n)+1+20n=40n$\\
今,$nが奇数なので,n\negodo 0\kmod{10}であるからf(n)\negodo0\kmod{400}である.$\\\\
$\Y 以上\tokeiichi,\tokeini より,f(n)を400で割り切るようなnは存在しない.\ans{(証明終わり)}$

\end{document}